\section{Results}
\label{sec:results}

\subsection{Dataset}

The campaign produced \num{63600} cases of \cA\ and \num{66400} of \cB, with no failed
simulation. All healthy cases are compliant. Of all cases, \SI{79.0}{\percent} of \cA\ and
\SI{69.8}{\percent} of \cB\ are compliant, although only \SI{7.9}{\percent} and \SI{7.5}{\percent}
are healthy: most injected faults do not take the circuit out of specification.
Fig.~\ref{fig:compliance} shows it for parametric faults. With one component \SI{5}{\percent} off
nominal, between 93 and \SI{96}{\percent} of the circuits remain compliant, and those that do not
are the ones in which the component sets the gain. With one component \SI{50}{\percent} off,
between 54 and \SI{66}{\percent} remain compliant. Among hard faults, \SI{46}{\percent} of the
open circuits and \SI{33}{\percent} of the short circuits of \cA\ leave it compliant (31 and
\SI{19}{\percent} in \cB).

\begin{figure}[!t]
\centering
\includegraphics[width=\columnwidth]{compliance_by_magnitude}
\caption{Circuits that still meet every specification when one passive component deviates from
its nominal value by the given percentage, the others staying within tolerance. Each point is the
mean over all the passive components of the circuit, 200 cases per component.}
\label{fig:compliance}
\end{figure}

In \cA\ no fault of the DRL integrator or of its input buffer (R7, R8, C4, U2, U3) violates a
specification, and neither do faults of the output and reference buffers (U5, U6); the passive
path through the output resistor of the DRL stage is enough for the \SI{89}{\decibel} test. Degrading the CMRR of the integrated amplifier to \SI{50}{\decibel} is
equally harmless, since the common-mode rejection of the system is dominated by the DRL loop.

The electrode has a large effect on what the self-test sees. The gain that a healthy \cA\
measures through $V_\mathrm{cal}$ at \SI{10}{\hertz}, relative to nominal, lies between 0.90 and
1.02 with gel and solid-metal electrodes, between 0.59 and 1.00 with the polymer, and between
0.05 and 0.97 with the fabrics (median 0.60 to 0.67). The megaohm-level contact impedance of
porous electrodes forms a divider with the input network. The circuit is nonetheless compliant,
because the standard tests the input impedance with \SI{620}{\kilo\ohm}.

\subsection{Testability}
\label{sec:res-testability}

\begin{table}[!t]
\caption{Testability of the Main Feature Set (Univariate Criterion)}
\label{tab:testability}
\centering
\setlength{\tabcolsep}{3pt}
\begin{tabular}{@{}p{0.50\columnwidth}cccc@{}}
\toprule
 & \multicolumn{2}{c}{\cA} & \multicolumn{2}{c}{\cB} \\
\cmidrule(lr){2-3}\cmidrule(l){4-5}
Electrodes & All & No porous & All & No porous \\
\midrule
Non-compliant conditions$^{\mathrm{a}}$ indistinguishable from healthy & 9 & 1 & 16 & 6 \\
Non-compliant conditions$^{\mathrm{a}}$ missed by the limit test & 15 & 3 & 35 & 8 \\
Testability rank (\SI{10}{\percent} deviation) & 2 & 4 & 2 & 3 \\
Components located without ambiguity & 4/23 & 5/23 & 2/29 & 5/29 \\
Mean number of components each one can be mistaken for & 3.9 & 2.5 & 8.6 & 5.2 \\
\bottomrule
\multicolumn{5}{@{}p{0.97\columnwidth}@{}}{\footnotesize $^{\mathrm{a}}$Fault conditions in
which most cases are non-compliant, out of 293 (\cA) and 307 (\cB).}
\end{tabular}
\end{table}

Table~\ref{tab:testability} summarizes the analysis of Section~\ref{sec:testability}. Three
observations follow.

First, some non-compliant faults cannot be seen. In \cA, eight of the nine invisible conditions
are gain errors of 5 to \SI{20}{\percent}, which are confused with the attenuation caused by a
high-impedance electrode; without porous electrodes only the open output resistor of the DRL
stage remains. In \cB\ six remain without porous electrodes, all in the DRL stage: they violate
the common-mode rejection requirement and the measurements do not see them.

Second, few components can be located without ambiguity, and the groups are predictable. The
sensitivities~\eqref{eq:sensitivity} alone give, for \cA, the groups R5+R6+R11+R12 (the resistors
that set the gain of the two stages), C5+R10 (high-pass filter), R13+R14 (low-pass filter) and
R15+R16 (reference divider). For \cB\ they give one group with the nine resistors that set the
gain (seven in the instrumentation amplifier and two in the gain stage), and the groups C4+R16
(high-pass filter), R19+R20 (low-pass filter) and R12+R13 (common-mode sensing).

Third, the structure depends on the architecture. \cA\ has fewer invisible faults and each
component can be mistaken for fewer others. The ratio of the distance between class centroids to
the spread within classes, a measure of the kind used in~\cite{chen2025}, is the only one that does not
favor it: 1.5 for both circuits with all electrodes, and 1.8 against 2.1 without porous
electrodes.

In the sensitivity analysis the common-mode tone responds to none of the components of either
circuit.

\subsection{Compliance Decision}
\label{sec:res-decision}

\begin{table}[!t]
\caption{Compliance Decision From the Main Feature Set. Mean of Three Repetitions and
\SI{95}{\percent} Confidence Interval, in Percent}
\label{tab:decision}
\centering
\setlength{\tabcolsep}{3.5pt}
\begin{tabular}{@{}llcc@{}}
\toprule
Circuit & Method & Escapes & False rejects \\
\midrule
A & Limit test & 22.5 (21.4--23.6) & 20.0 (19.8--20.3) \\
 & Specification regression & 8.4 (6.7--10.2) & 1.5 (1.3--1.6) \\
 & Classifier & 7.9 (6.4--9.4) & 0.7 (0.4--1.0) \\
 & Classifier, \SI{1}{\percent} escape target & 0.9 (0.8--1.0) & 7.8 (4.8--10.8) \\
\addlinespace
B & Limit test & 35.3 (34.7--35.8) & 17.0 (16.5--17.6) \\
 & Specification regression & 11.4 (10.0--12.8) & 2.4 (1.9--2.9) \\
 & Classifier & 11.2 (10.3--12.1) & 1.2 (1.1--1.2) \\
 & Classifier, \SI{1}{\percent} escape target & 1.2 (0.6--1.7) & 54.0 (48.0--60.0) \\
\addlinespace
\multicolumn{4}{@{}l}{\emph{Without porous dry electrodes}} \\
A & Classifier & 4.2 (3.5--4.9) & 0.7 (0.4--0.9) \\
 & Classifier, \SI{1}{\percent} escape target & 1.0 (0.9--1.1) & 1.9 (1.3--2.5) \\
B & Classifier & 6.7 (6.2--7.1) & 1.0 (0.6--1.4) \\
 & Classifier, \SI{1}{\percent} escape target & 1.0 (0.7--1.3) & 48.7 (39.4--58.1) \\
\bottomrule
\end{tabular}
\end{table}

\begin{figure*}[!t]
\centering
\includegraphics[width=\textwidth]{operating_points}
\caption{Operating points of the four decision methods on the main feature set, for all electrode
types (filled markers) and without porous dry electrodes (open markers). Bars are \SI{95}{\percent}
confidence intervals over three repetitions. Both axes are logarithmic.}
\label{fig:operating}
\end{figure*}

Table~\ref{tab:decision} and Fig.~\ref{fig:operating} give the result. The limit test is poor in
both senses: it accepts between a fifth and a third of the non-compliant circuits and rejects a
fifth of the compliant ones, because a wide healthy envelope (widened further by the electrodes)
and a narrow compliance region do not coincide. The specification regression and the classifier
are close to each other at their default operating points. Per specification, the error of the
regression is between 8 and \SI{36}{\percent} of the limit, the largest for the input-impedance
test.

An instrument needs a low escape rate. Lowering the threshold of the classifier brings escapes
on \cA\ to \SI{0.9}{\percent} at the cost of \SI{7.8}{\percent} of false rejects, and of only
\SI{1.9}{\percent} without porous dry electrodes. On \cB\ the same escape rate costs more than
half of the compliant circuits. The escapes that remain there are the DRL faults that the
measurements do not see, as the testability analysis anticipated; no threshold can recover a
fault that leaves no trace. The guard band on the regression does not provide a usable operating
point: reaching \SI{1}{\percent} of escapes rejects more than \SI{90}{\percent} of the compliant
circuits in both designs.

Telling the classifier which electrode type is connected does not help (\SI{8.3}{\percent} of
false rejects at the same escape target on \cA). The variation that matters is between subjects,
not between materials.

No single set of measurements is sufficient. With the escape target, the pulse response alone
(C3) gives \SI{44}{\percent} of false rejects on \cA, and each of the other sets alone gives more
than \SI{75}{\percent}.

\subsection{Functional Against Percentage Severity}
\label{sec:res-severity}

\begin{table}[!t]
\caption{Detector Trained on the Percentage Label or on the Functional Label, Scored Against
Compliance. Mean of Three Repetitions, in Percent}
\label{tab:severity}
\centering
\begin{tabular}{@{}llccc@{}}
\toprule
Circuit & Trained on & Labeled bad & Escapes & False rejects \\
\midrule
A & Percentage & 88.7 & 7.9 & 36.8 \\
 & Functional & 21.0 & 4.1 & 1.6 \\
B & Percentage & 89.2 & 8.3 & 36.5 \\
 & Functional & 30.2 & 7.6 & 2.2 \\
\bottomrule
\end{tabular}
\end{table}

Table~\ref{tab:severity} compares the two definitions of a fault. A detector trained to recognize
that a component is out of tolerance rejects \SI{37}{\percent} of the circuits that comply with
the standard, in both designs. Trained on compliance, the same model rejects \SI{1.6}{\percent}
and \SI{2.2}{\percent}, and it also has fewer escapes. The false rejects of the percentage-trained
detector are dominated by parametric faults that are real but harmless. The escape rate of the
functional detector is lower here than in Table~\ref{tab:decision} because this experiment
weights the classes, which moves the operating point.

\subsection{Location}
\label{sec:res-location}

\begin{table}[!t]
\caption{Location of the Cause on Non-Compliant Cases: Macro F1 (Mean of Three Repetitions)}
\label{tab:location}
\centering
\begin{tabular}{@{}lcccc@{}}
\toprule
 & \multicolumn{2}{c}{By ambiguity group} & \multicolumn{2}{c}{By component} \\
\cmidrule(lr){2-3}\cmidrule(l){4-5}
Model & A & B & A & B \\
\midrule
$k$-nearest neighbors & 0.83 & 0.70 & 0.67 & 0.55 \\
Random forest & 0.86 & 0.76 & 0.74 & 0.61 \\
Support vector machine & 0.86 & 0.71 & 0.69 & 0.53 \\
Gradient boosting & 0.87 & 0.74 & 0.73 & 0.59 \\
Multilayer perceptron & 0.81 & 0.71 & 0.66 & 0.53 \\
Convolutional network$^{\mathrm{a}}$ & 0.56 & 0.32 & -- & -- \\
\bottomrule
\multicolumn{5}{@{}l}{\footnotesize $^{\mathrm{a}}$On the sampled pulse response only.}
\end{tabular}
\end{table}

\begin{figure}[!t]
\centering
\includegraphics[width=\columnwidth]{recall_by_component}
\caption{Recall of each component by the best component-level model (random forest) on
non-compliant cases. Hatched bars are members of the ambiguity groups predicted from the
sensitivities of the nominal circuit; a bracket joins the members of each group.}
\label{fig:recall}
\end{figure}

There are \num{13364} non-compliant cases of \cA\ over 23 components and \num{20033} of \cB\ over
29. Table~\ref{tab:location} gives the macro F1 of each model. At the level of the a priori
ambiguity groups the best models reach 0.87 on \cA\ and 0.76 on \cB; with the random forest the
classification accuracy is 95 and \SI{94}{\percent}, and the true group is among the three most probable in
99.9 and \SI{99.4}{\percent} of the cases. At component level the macro F1 falls to 0.74 and
0.61.

Fig.~\ref{fig:recall} shows where the component-level errors are. In \cA\ they concentrate in the
predicted groups: the two halves of the gain resistor (R5, R6), the two resistors of the low-pass
filter (R13, R14) and, to a lesser degree, the gain-stage and reference resistors. The recall of a
component decreases with the number of components that the testability analysis predicted it
could be mistaken for (Spearman $\rho=-0.42$, $p=0.045$ for \cA; $\rho=-0.51$, $p=0.005$ for
\cB). In \cB\ the nine resistors of the gain group are all located with a recall between 29 and
\SI{64}{\percent}.

Two components of \cA\ outside the predicted groups are also located poorly: the bypass capacitor
of the reference (C8), which the analysis of simulated faults, though not the sensitivities,
links to the reference divider, and the integrated amplifier (U1), whose non-compliant faults are
gain errors indistinguishable from those of its gain resistors. The convolutional network on the raw pulse response is clearly below
the tabular models that use all the measurements.

\subsection{Electrode or Circuit}
\label{sec:res-origin}

\begin{table}[!t]
\caption{Origin Classification on the Test Part of One Repetition (Random Forest, Main Feature
Set)}
\label{tab:origin}
\centering
\begin{tabular}{@{}llcccc@{}}
\toprule
 & & \multicolumn{3}{c}{Predicted} & \\
\cmidrule(lr){3-5}
Circuit & True & None & Circuit & Electrode & Recall \\
\midrule
A & None & \num{14314} & 63 & 60 & \SI{99.2}{\percent} \\
 & Circuit & 189 & \num{3649} & 25 & \SI{94.5}{\percent} \\
 & Electrode & 269 & 15 & 496 & \SI{63.6}{\percent} \\
\addlinespace
B & None & \num{13029} & 136 & 68 & \SI{98.5}{\percent} \\
 & Circuit & 445 & \num{5417} & 45 & \SI{91.7}{\percent} \\
 & Electrode & 278 & 25 & 477 & \SI{61.2}{\percent} \\
\bottomrule
\end{tabular}
\end{table}

Table~\ref{tab:origin} shows the three-class result; the macro F1 over three repetitions is 0.89
for \cA\ and 0.87 for \cB. A case with nothing to act on is recognized in \SI{99}{\percent} of
the cases and a non-compliant circuit in 94 and \SI{92}{\percent}. Circuit faults are almost
never taken for electrode problems.

Electrode problems are recognized in only 64 and \SI{61}{\percent} of the cases, and the
breakdown by type explains why. An open protection resistor is always recognized. A detached
electrode is recognized in 75 and \SI{89}{\percent} of the cases. A degraded contact is
recognized in half of the cases in \cA\ and in \SI{41}{\percent} in \cB, and otherwise passes as
normal: a contact impedance five or twenty times higher than that of the same subject falls
within what is measured on other subjects with healthy contacts. Without porous dry electrodes
the recall of the electrode class rises to 76 and \SI{74}{\percent}.

\subsection{Robustness}
\label{sec:res-robustness}

\begin{figure*}[!t]
\centering
\includegraphics[width=\textwidth]{unseen_magnitudes}
\caption{Location of the cause for parametric faults of magnitudes that were present in training
or held out of it entirely, scored on the non-compliant cases of those magnitudes. The model is
trained to name the component; the group bars count a prediction as correct when it falls in the
ambiguity group of the true component.}
\label{fig:unseen}
\end{figure*}

\subsubsection{Unseen magnitudes}
The compliance decision generalizes to fault magnitudes that were not in the training data: its
balanced accuracy on the held-out magnitudes falls by less than 2 points with respect to a model
trained with every magnitude, except for \SI{50}{\percent} deviations on \cA\ (6 points). Component-level location does not (Fig.~\ref{fig:unseen}): the
share of correctly located components falls from between 21 and \SI{70}{\percent} to between 2 and
\SI{33}{\percent}. Within a group of components with parallel sensitivities the model can only
tell the members apart by the size of the effect, that is, by memorizing the magnitudes it was
shown. Counted at group level, the same predictions remain correct in 99 and \SI{95}{\percent} of
the cases on \cA\ for held-out small and intermediate magnitudes and in \SI{84}{\percent} for
large ones (96, 81 and \SI{79}{\percent} on \cB).

\subsubsection{Noise, resolution and tolerances}
Table~\ref{tab:robustness} gives the remaining results. When the models are trained at the same
noise level and resolution at which they are used, the results barely change from 0.1 to
\SI{10}{\milli\volt} of noise and from 8 to 16 bits. When they are trained at
\SI{1}{\milli\volt} and used at \SI{10}{\milli\volt}, escapes on \cA\ rise from 4 to
\SI{10}{\percent} and false rejects on \cB\ from 2 to \SI{19}{\percent}. A different tolerance
distribution changes nothing. Looser components (\SI{2}{\percent} resistors and \SI{10}{\percent}
capacitors) raise escapes in both circuits; group-level location is unaffected in \cA\ and
degrades strongly in \cB.

\begin{table}[!t]
\caption{Robustness of the Compliance Decision (Class-Weighted Classifier) and of Group-Level
Location. Mean of Three Repetitions}
\label{tab:robustness}
\centering
\setlength{\tabcolsep}{3pt}
\begin{tabular}{@{}p{0.40\columnwidth}cccccc@{}}
\toprule
 & \multicolumn{3}{c}{\cA} & \multicolumn{3}{c}{\cB} \\
\cmidrule(lr){2-4}\cmidrule(l){5-7}
Condition & Esc. & F.r. & F1 & Esc. & F.r. & F1 \\
\midrule
Default (\SI{1}{\milli\volt}, 12 bits) & 4.1 & 1.7 & 0.88 & 7.6 & 2.2 & 0.77 \\
\SI{10}{\milli\volt} noise, matched training & 4.2 & 1.6 & 0.86 & 7.8 & 2.2 & 0.75 \\
\SI{10}{\milli\volt} noise, trained at \SI{1}{\milli\volt} & 9.7 & 1.2 & 0.79 & 6.6 & 18.8 & 0.71 \\
8 bits, matched training & 3.9 & 1.6 & 0.88 & 7.5 & 2.2 & 0.76 \\
8 bits, trained at 12 bits & 7.0 & 1.2 & 0.87 & 7.2 & 4.2 & 0.75 \\
Truncated Gaussian tolerances & 3.8 & 1.6 & 0.89 & 7.4 & 2.4 & 0.77 \\
\SI{2}{\percent} resistors, \SI{10}{\percent} capacitors & 7.1 & 1.9 & 0.87 & 11.4 & 2.9 & 0.45 \\
\bottomrule
\multicolumn{7}{@{}p{0.97\columnwidth}@{}}{\footnotesize Esc.: escapes (\%). F.r.: false rejects
(\%). F1: macro F1 of the location by ambiguity group. The last two rows are scored on
separately generated datasets.}
\end{tabular}
\end{table}

\subsubsection{Common-mode tone}
Raising the amplitude of the common-mode tone from 0.1 to \SI{1}{\volt}, or making it ten times
longer, does not make the faults of the DRL stage detectable: about half of the non-compliant
cases of that stage escape whatever the setting (51 to \SI{45}{\percent} in \cA, 53 to
\SI{55}{\percent} in \cB). The limitation is not one of signal-to-noise ratio; a different
measurement is needed.

\subsection{Minimal Set of Measurements}
\label{sec:res-minimal}

\begin{figure*}[!t]
\centering
\includegraphics[width=\textwidth]{cost_performance}
\caption{Score on the test part as self-test measurements are added one at a time by greedy
selection on the validation part, against the accumulated duration of the self-test. One line
per repetition. Dashed lines are the score with all the measurements.}
\label{fig:cost}
\end{figure*}

Fig.~\ref{fig:cost} shows the score as measurements are added; the candidates are the 15 actions
of Table~\ref{tab:measurements} and the dc reading of the DRL output. The calibration pulse is
always selected first. For the compliance decision on \cA, three measurements taking \SI{3}{\second}
(the pulse, one lead-off tone and the differential tone at \SI{150}{\hertz}) are within one point
of the complete self-test of \SI{97}{\second}, in the three repetitions. \cB\ needs a fourth, the
dc reading of the DRL output, which is the measurement that sees its DRL faults. For location on
\cA, the pulse, the \SI{150}{\hertz} tone and the dc reading of the instrumentation-amplifier
output (\SI{2}{\second}) reach the score of the complete set in two of the three repetitions; \cB\
needs five or six measurements. The \SI{0.05}{\hertz} tones, which account for 80 of the
\SI{97}{\second}, are never among the first three measurements selected.
